\section{Original NOTEARS: preserving the numerical meaning of a discovery reference}
\label{exec23:notears-scope}
A bank investigating funding transmission needs to know whether a proposed
edge represents a signed structural coefficient, a measure of connectivity,
or a decision made by thresholding. The original \texttt{xunzheng/notears}
implementation and the inspected gCastle implementation are separate software
candidates. Their source revisions, threshold operators and numerical
constraints must remain visible in a comparison. The original source is
retained here as a reference for a proposed public prototype; its output does
not establish bank identification or calibrated effect uncertainty.
\citep{exec23_notears_source,exec20_gcastle_source}

The selected original revision precedes the 7 October 2026 cutoff. Its
setup declares version 3.0, an alpha classifier and Apache License 2.0,
with unpinned NumPy, SciPy and python-igraph requirements. The nonlinear
path also imports PyTorch, although it is absent from the setup requirements;
selected Gaussian-process simulation paths import scikit-learn. The root
license text and optimizer wrapper attribution are retained in the source
audit. These observations establish the selected source and declared terms;
they do not establish a resolved distribution or transitive dependency review.
The README points to DAGrad for newer differentiable methods. That pointer is
a separate candidate lead, not evidence of a validated successor or a measured
performance advantage.\citep{exec23_notears_source}

\subsection{A weighted matrix is a contract, not a complete causal result}
\label{exec23:notears-matrix}
For an $n\times d$ matrix $X$ of observations, the original linear routine uses
$XW$ to predict columns. Thus $W_{ij}$ is the coefficient associated with
source variable $i$ and target variable $j$. The least-squares path centers a
new data array; logistic and Poisson losses use different response models.
A wrapper must preserve the selected loss, variable order, units, centering
and sampling interpretation. The routine consumes returned optimizer
parameters without exposing its success status. Termination can follow an
acyclicity tolerance, a penalty bound or the iteration budget. Its public
function returns a thresholded matrix. A graph assertion appears in the
standalone demonstration and CLI, but not as a final guard inside the linear
function. An adapter must check the returned support itself and retain a
reason when it cannot certify a DAG.\citep{exec23_notears_source}

The original CLI reads a numerical CSV, invokes the routine, asserts
acyclicity and saves a matrix. It does not serialize variable identities,
solver termination, the information clock or uncertainty semantics. Those
fields belong in the proposed common library, CLI and MCP operation contract;
they cannot be reconstructed from matrix dimensions after a restart. Bounds on
elapsed time, nonfinite results and retries must also contribute to the failure
denominator. A missing output cannot disappear from the comparison.
\citep{exec23_notears_source}

\subsection{Threshold equality can change a banking decision}
\label{exec23:notears-threshold}
The original removes an entry when $|W_{ij}|<\tau$, retaining equality.
The inspected gCastle binary output keeps an entry when $|W_{ij}|>\tau$;
its weighted output remains a separate object. Even on a supplied identical
matrix, the two support conventions disagree at equality. This is a conversion
boundary, not an empirical ranking of fitted algorithms.
\citep{exec23_notears_source,exec20_gcastle_source}

Consider a stipulated linear causal relation between funding support and a
monetary response, with coefficient $3/10$ dollars per intervention unit.
Let the contemplated intervention be 100 units and its cost USD 20.
At a threshold $\tau=3/10$, retaining the supplied coefficient gives
\begin{equation}
 \Delta V=100(3/10)-20=10\quad\text{USD}.
 \label{exec23:notears-value}
\end{equation}
If a consumer reconstructs a coefficient model from a binary support that
removes this boundary edge, its implied response is zero and its computed
margin becomes USD $-20$. This deliberately constructed discrepancy concerns
using a support mask consistently; it does not say that either package has
estimated the coefficient or that graph selection justifies a causal effect.
The source-row convention, threshold strictness and distinction between
weights and support must accompany the proposed output.
\citep{exec23_notears_reference}

Units add another constraint. With $Z=XD^{-1}$ for a positive diagonal scaling
matrix $D$, the same supplied linear equations have $W_Z=DW_XD^{-1}$.
For the three-variable chain with $W_{12}=2$, $W_{23}=3$ and $W_{13}=1$,
the total path coefficient is $1+2\cdot3=7$.
If $D=\operatorname{diag}(10,1,100)$, the transformed coefficient is $7/10$.
The change is exactly the change in units. A fixed numerical threshold need
not preserve support under this transformation. This is an algebraic conversion
identity; refitting a penalized objective after scaling need not reproduce
the transformed original fit.\citep{exec23_notears_reference}

\subsection{A small residual need not certify a DAG}
\label{exec23:notears-residual}
For the supplied two-variable matrix
$W=\left(\begin{smallmatrix}0&a\\b&0\end{smallmatrix}\right)$,
the original exponential residual is
\begin{equation}
 h_{\exp}(W)=\operatorname{tr}\exp(W\circ W)-2
            =2\{\cosh(ab)-1\}.
 \label{exec23:notears-cycle}
\end{equation}
Here $\circ$ denotes elementwise multiplication. If $a=b=1/200$,
then $z=ab=1/40000$. Expanding the hyperbolic cosine gives exact bounds
\begin{equation}
 z^2<h_{\exp}(W)\leq\frac{z^2}{1-z^2/12}<10^{-8}.
 \label{exec23:notears-bound}
\end{equation}
The upper bound follows because successive positive terms after $z^2$
have ratios at most $z^2/12$. Both edges nevertheless remain at a deliberately
chosen threshold $1/1000$. The support is a directed cycle. The default
threshold $3/10$ would remove these edges; no actual optimizer output is
asserted by this construction. It distinguishes a numerical stopping
tolerance from an exact graph check.
\citep{exec23_notears_source,exec23_notears_reference}

The polynomial trace residual used in the selected gCastle linear routine
is a different numerical quantity. At $d=2$ it equals $a^2b^2/2$.
For $a=b=1/100$ it is $1/(2\cdot10^8)$, below $10^{-8}$,
while the exponential residual is strictly greater than $10^{-8}$.
Passing the same tolerance under the two formulas is not an identical
numerical event. A matched study should report the actual formula, value,
final support and stopping reason rather than compare an unqualified
convergence flag.\citep{exec20_gcastle_source,exec23_notears_reference}

\subsection{Nonlinear adjacency cannot stand in for a signed response}
\label{exec23:notears-nonlinear}
The original MLP and Sobolev paths return nonnegative norms of first-layer
weights. Such an adjacency summarizes connectivity and cannot recover the
sign or magnitude of an intervention response. The MLP uses local output
layers after the first layer; the model is mutated by optimization. Its
SciPy bridge moves through NumPy and the default PyTorch dtype, and the
trace-exponential helper crosses a NumPy/SciPy boundary. Device and dtype
compatibility are obligations for a future installed prototype, not inferred
from the existence of a PyTorch class.\citep{exec23_notears_source}

Let $s(x)=1/(1+e^{-x})$ and consider two stipulated structural functions
$f_+(x)=4s(2x)-2$ and $f_-(x)=-4s(2x)+2$.
Both have first-layer norm 2 and baseline output zero. Since $s'(0)=1/4$,
their local derivatives at zero are respectively 2 and $-2$.
With assumed value USD 100 per outcome unit and local cost USD 20 per
intervention-intensity unit, the net local derivatives are USD 180 and
USD $-220$. These are local derivatives in stipulated models, not finite
policy returns. Identical adjacency norms cannot choose between them.
A downstream intervention query needs the fitted structural functions,
noise assumptions, unit conversion and an identification argument.
\citep{exec23_notears_reference}

\subsection{Reuse choice and the proposed public experiment}
\label{exec23:notears-choice}
A pinned original linear worker is a reasonable prototype route because it
preserves the reference method while a wrapper owns input validation, source
identity, diagnostics and graph semantics. Direct adoption of the entire
package is deferred: the declared environment is not locked, nonlinear outputs
have different meaning, and solver status is not exposed by the inspected
entrypoints. A port into another language is not selected for this proposal
because there is no demonstrated deployment or cost benefit to offset new
numerical and maintenance obligations. gCastle and DAGMA remain separately
reviewed comparators; they do not inherit one another's qualification.
\citep{exec23_notears_source,exec20_gcastle_source,exec16_dagma_primary}

The selected sources warn against treating bundled experiments as completed
reproducibility evidence. The two-variable experiment retains a legacy import
and function name absent from the selected package tree. The nonlinear CLI
constructs a randomly initialized model without setting an explicit seed;
the utility seed function sets Python and NumPy state, while local layer
initialization uses PyTorch. These are static observations requiring native
compatibility and repeatability checks. No downloaded experiment was run
for this review.\citep{exec23_notears_source}

A prospective comparison should first replay the conversion examples, then
use predeclared public or generated mechanisms under matched causal and
sampling restrictions. It must count refused, failed and incomplete fits;
retain variable permutations, scale changes, near-threshold cases and weak
signals; and distinguish optimizer, topology and statistical diagnostics.
The scientific owner must justify model restrictions and effect queries; the
maintainer must own the environment and failure contract; the future operator
must resolve actual distribution terms and custody. Installation, full job
cost, serialization, diagnostics and upgrade effort should be measured before
arguing for costly bank experiments. The reference decision supplies that
experimental rationale without declaring an implemented or qualified bank
product.\citep{exec23_notears_source,exec23_notears_reference}
